% E6 — hazard-guided adaptive controller. Slots into results_analysis.tex as the
% closing subsection (\S\ref{sec:rq6}). Figures: results/figures/e6_controller.png,
% results/figures/e6_collateral.png. Source table: results/tables/e6_controller.json.
%
% Terminology discipline (do NOT revert):
%   * "hazard-guided controller" — NOT "influence-guided". E2 showed influence alone
%     is insufficient to attribute the mechanism; at the g=0 gate the detector uses
%     round-0 disagreement geometry (F1) + text (F4) + influence (F3). F2 dynamics are
%     ZERO at g=0 (nothing has revised yet) -- do not claim "revision dynamics" here.
%   * tau* is an OBSERVED Pareto point, not a tuned knob (E12 cross-fits it).
%   * NUMBERS ARE PER-SEED (fixed cache, keyed bench:qid:seed). The earlier per-qid
%     cache collapsed the 5 seeds -> inflated headroom (8.5->4.5pp) and rescue
%     (.46->.40) and a spurious gsm8k advantage. Do NOT reintroduce old numbers.
%   * Honest headline (E9, cluster-bootstrap over 184 questions): the controller
%     MATCHES always-on (hazard-always = +0.0pp) at ~55% fewer interventions; its edge
%     over random is borderline (+2pp pooled, ~0 within each benchmark); the ONE
%     significant result is the rescue>>harm asymmetry (0.40 vs 0.21).

\subsection{Selective intervention is efficient: a hazard-guided controller}
\label{sec:rq6}

The preceding results close a chain of inference. Debate can \emph{destroy} a correct
answer that was alive early (\S\ref{sec:rq4}, E1); this harm has a reproducible
anatomy (E4a); the round-0 dynamics make it \emph{predictable} before it happens
(E3a); and a targeted re-query \emph{repairs} it when it is about to occur (E5). The
open question is whether prediction and repair can be combined into a policy that is
better than intervening everywhere. Crucially, the triggering signal is \emph{not}
influence: E2 established that influence alone does not attribute the failure
mechanism. The controller is therefore \emph{hazard-guided} --- it fires on the E3a
Stage-A hazard, a leave-one-group-out out-of-fold score built from the round-0
behavioral state: answer-disagreement geometry (F1), reasoning-text features (F4), and
global influence (F3). Revision dynamics (F2), which strengthen the pooled detector at
later observation times, are by construction unavailable this early --- the earliest
deaths occur before any agent has revised. \S\ref{sec:rq7} shows that \emph{no single
channel dominates}: structure, text, and influence each contribute, influence is
neither necessary nor sufficient, and the combined detector is the robust choice.

\paragraph{Setup.}
We evaluate on the $n=516$ debates ($184$ distinct questions $\times$ up to five seeds)
that have a correct answer alive at round~0 --- the population that \emph{has} something
to protect (extinction prevalence $22.5\%$). For each debate we precompute, \emph{per
seed}, two counterfactual continuations from a single early gate: \textsc{control} (a
plain stochastic re-query, the E5 baseline) and \textsc{combined} (self-verification
plus an influence-resistant re-query). Every controller is then an \emph{offline}
selection policy over these two stored outcomes, so the entire accuracy--intervention
curve is swept at zero additional GPU cost. Two anchors bound the curve: a \emph{never}
policy (\textsc{control} everywhere; accuracy $58.3\%$) and an \emph{always-on} policy
(\textsc{combined} everywhere; accuracy $62.8\%$), a headroom of $4.5$ points. The
\emph{hazard-guided} policy
applies \textsc{combined} iff the E3a hazard exceeds a threshold $\tau$; a
\emph{random-matched} policy spends the identical intervention budget on a random
subset and is our null model for ``does the detector add anything.''

\paragraph{Selective beats indiscriminate.}
Figure~\ref{fig:e6}A traces accuracy against intervention rate. In point estimate the
hazard-guided controller tracks above the random-matched budget across the sweep and
\emph{numerically exceeds} always-on at a much reduced intervention rate: at its
observed peak operating point ($\tau^\star=0.52$) it reaches
\[
  \mathrm{Acc}_{\text{hazard}} = 63.2\%
  \;\gtrsim\;
  \mathrm{Acc}_{\text{always}} = 62.8\%
  \;>\;
  \mathrm{Acc}_{\text{never}} = 58.3\%,
\]
while intervening on only $45\%$ of debates ($55\%$ fewer safeguard
\emph{interventions} than always-on). Under a cluster-bootstrap over the $184$
independent questions (\S\ref{sec:rq9}), the hazard$-$always margin is $+0.0$ points
(CI $[-1.6,+1.2]$), squarely within noise, so the defensible claim is that the
controller \emph{matches} always-on accuracy at roughly half the interventions; its
edge over a matched-random budget is \emph{suggestive but not conclusive} ($+2$ points
pooled, near zero within either benchmark). The one comparison that \emph{is}
significant is the collateral asymmetry below. This is not a compute
argument: the combined safeguard adds only $\sim\!0.06$ model calls per debate on
average ($5.65\!\to\!5.72$), so raw token cost is essentially unchanged; the
controller's value is doing the job of always-on with far fewer interventions and less
collateral harm. Figure~\ref{fig:e6}B recasts this as
gain retention, $(\mathrm{Acc}-\mathrm{never})/(\mathrm{always}-\mathrm{never})$.
Because the always$-$never headroom is only $4.5$ points, gain retention is a
high-variance summary; we therefore rely on the raw accuracy edge over the
matched-random budget ($\approx\!+2$ points pooled at a $40\%$ intervention rate) with
its cluster-bootstrap interval (\S\ref{sec:rq9}) and on the baseline comparisons
(\S\ref{sec:rq7}), rather than on any single gain-retention number.

We present $\tau^\star$ as an \emph{observed} Pareto operating point, not a tuned
hyperparameter, and \S\ref{sec:rq12} confirms by cross-fitting that a data-chosen
threshold does not inflate the result (held-out accuracy matches always-on). The
conclusion rests on the full budget frontier, not any single threshold.

\paragraph{Why a controller has to exist.}
The gain is not free efficiency --- it reflects that intervention is
\emph{asymmetric} (Figure~\ref{fig:e6collateral}). Applying the combined safeguard
indiscriminately rescues $39.5\%$ of would-die debates ($n=215$, those the plain
re-query leaves wrong) but \emph{harms} $20.6\%$ of would-live debates ($n=301$,
those that were already correct and are knocked off the answer by the extra
perturbation). Always-on intervention pays this collateral tax on every safe debate;
the hazard-guided controller exists precisely to concentrate intervention where its
expected value is positive and to withhold it where the dominant risk is
self-inflicted harm. This asymmetry (rescue $\gg$ harm) is the one result that
\emph{is} statistically significant under question-level clustering (\S\ref{sec:rq9}),
and it is the controller's principled justification: even where selective and
indiscriminate intervention tie on aggregate accuracy, targeting is preferable because
it avoids the collateral damage that blanket safeguarding necessarily incurs.

% ---------------------------------------------------------------------------
\begin{figure}[t]
  \centering
  \includegraphics[width=\linewidth]{results/figures/e6_controller.png}
  \caption{\textbf{A hazard-guided controller makes selective intervention efficient}
  ($n=516$ debates with a correct answer alive at round~0).
  \textbf{(A)} Debate accuracy vs.\ intervention rate. In point estimate the
  hazard-guided controller (blue) tracks above a random-matched budget (grey, within-run
  $95\%$ CI) and \emph{numerically exceeds} always-on: at the observed operating point
  $\tau^\star=0.52$ (star) it reaches $63.2\%$ vs.\ the always-on $62.8\%$, while
  intervening on only $45\%$ of debates (\emph{never} floor $58.3\%$). Under a proper
  cluster-bootstrap over questions (\S\ref{sec:rq9}) the hazard$-$always margin is $+0.0$
  points (within noise) and the hazard$-$random edge is borderline; the robustly
  significant result is the collateral asymmetry (Fig.~\ref{fig:e6collateral}).
  \textbf{(B)} The same sweep as gain retention. In the $20$--$40\%$ operating band
  (shaded) the controller clears the $\mathrm{GR}=0.70$ gate in point estimate and lies
  above the matched-random budget; its vertical gap over the grey curve is the
  detector's (point-estimate) value. $\tau^\star$ is an observed Pareto point, not a
  tuned threshold (cross-fitted in \S\ref{sec:rq12}).}
  \label{fig:e6}
\end{figure}

\begin{figure}[t]
  \centering
  \includegraphics[width=0.72\linewidth]{results/figures/e6_collateral.png}
  \caption{\textbf{Intervention is asymmetric --- the controller's reason to exist.}
  Applying the combined safeguard everywhere rescues $39.5\%$ of would-die debates
  ($n=215$) but inflicts collateral harm on $20.6\%$ of would-live debates
  ($n=301$); the two intervals do not overlap (\S\ref{sec:rq9}), so the asymmetry is
  significant. Because indiscriminate safeguarding taxes every already-correct debate,
  a hazard-guided policy that withholds intervention on low-hazard debates does the job
  of always-on with far fewer interventions.}
  \label{fig:e6collateral}
\end{figure}
